% page 253 du fac-simile (image p-252 du scan)
% bloc 165.1 x 172.8 mm ; marge gauche 36.9 mm
\pgc{15.240mm}{0.000mm}{
\l{\cel{52}245}
\l{}
\l{}
\l{}
\l{}
\l{}
\l{\sou{izoklino}.\cel{2}(\sou{magnetismo}). Lineo ek la punti di la Terglobo ube}
\l{\cel{2}la deviaco da la magnet-indexo esas intersama. - DEFIS.}
\l{}
\l{\sou{izokron}a.(\sou{fiziko}). Qua eventas sama-dure. - DEFIS.}
\l{}
\l{\sou{izola}r. (\sou{trans.,}\cel{1}de)\cel{2}I.Separar de omna kozo : (a) igar (korpo)}
\l{\cel{2}ne plus kontaktar kun korpi elektro-konduktiva ; (b) separar}
\l{\cel{2}korpi ek la elementi kemiala kun qui lu esas kombinita ; (c)}
\l{\cel{2}igar (ulu) vivar for lua kunhomi. - DEFIRS.}
\l{}
\l{\sou{izomer}a.(\sou{kemio)}. Qua konsistas ye atomi qualeso-sama e quanteso-}
\l{\cel{2}sama. - DEFIS.}
\l{}
\l{\sou{izometr}a.(\sou{geom.})\cel{2}(\sou{pri perspektivo)}.\cel{2}Olta en qua la axi kompa-}
\l{\cel{2}rala esas inter-egala. - DEFIRS.}
\l{}
\l{\sou{izomorf}a.(\sou{kemio}).\cel{2}Di qua la formo kristala esas sama kam olta}
\l{\cel{2}di korpo kemiala altra.- DEFIRS.}
\l{}
\l{\sou{izoperimetr}a. (\sou{geom}.) . Dicesas pri la figuri di qui la peri-}
\l{\cel{2}metri esas inter-egala.}
\l{}
\l{\sou{izoterm}o.(\sou{fiziko)}.\cel{2}Lineo, sur mapo, qua pasas tra omna punti}
\l{\cel{2}di la Terglobo ube la temperaturo mezvalora dum la yaro esas}
\l{\cel{2}inter-sama. - DEFIRS.}
\l{}
\l{\sou{izotonik}a. (\sou{biol.})\cel{2}Dicesas pri la equilibro molekulala di du}
\l{\cel{2}solvuri interseparita da membrano organika, e qui havas}
\l{\cel{2}osmo-povo interegala.}
\l{}
\l{\sou{izotop}a.(\sou{kemio}). Dicesas pri la kompozaji kemiala inter-identa,}
\l{\cel{2}ma di qui la pezo atomala interdiferas. - F.}
\l{}
\l{\sou{izotrop}a. Qua havas la sama partikularaji fizikala omna-direcione}
}
